\uniterablesection{Вступ}

Даний проект присвячений розробці автоматизованого робочого місця (АРМ) адміністратора станції технічного обслуговування. Впровадження такої системи дозволяє централізувати управління замовленнями, забезпечити прозорий облік ресурсів та мінімізувати помилки при оформленні документації.

\subsection{Постановка проблеми}
Основним викликом для власника або адміністратора СТО є необхідність одночасного контролю великої кількості параметрів: від реєстрації клієнта до фінального розрахунку вартості. Без спеціалізованого програмного забезпечення адміністратору важко:
\begin{itemize}
    \item Вести єдиний реєстр замовлень без дублювання даних.
    \item Оперативно відстежувати, який майстер вільний, а який завантажений.
    \item Точно розраховувати суму до оплати з урахуванням використаних запчастин зі складу.
\end{itemize}

\subsection{Мета та функціонал системи}
Метою роботи є створення зручного інтерфейсу для адміністратора, який забезпечує повний цикл супроводу клієнта. 

Система дозволяє адміністратору виконувати наступні дії:
\begin{itemize}
    \item \textbf{Управління заявками:} Створення нових карток ремонту з описом несправностей та даними автомобіля.
    \item \textbf{Диспетчеризація:} Призначення конкретного виконавця (майстра) на кожне замовлення.
    \item \textbf{Складський облік:} Списання запчастин, використаних у процесі ремонту, безпосередньо через форму замовлення.
    \item \textbf{Фінансовий контроль:} Автоматичне формування підсумкової вартості замовлення за формулою: $V_{total} = V_{work} + V_{parts}$
\end{itemize}

\subsection{Об'єктно-реляційна модель}
Для реалізації функцій адміністратора база даних проекту містить взаємопов'язані таблиці:
\begin{itemize}
    \item \textbf{Клієнти та Автомобілі} — база контактів та об'єктів ремонту.
    \item \textbf{Співробітники} — реєстр майстрів, якими оперує адміністратор.
    \item \textbf{Замовлення} — центральна таблиця, що об'єднує клієнта, майстра та статус виконання.
\end{itemize}
